%Daten zur Institution

%\input{Dozentdaten}

%\renewcommand{\fachbereich}{Fachbereich}

%\renewcommand{\dozent}{Prof. Dr. . }

%Klausurdaten

\renewcommand{\klausurgebiet}{ }

\renewcommand{\klausurtyp}{ }

\renewcommand{\klausurdatum}{ . 20}

\klausurvorspann {\fachbereich} {\klausurdatum} {\dozent} {\klausurgebiet} {\klausurtyp}

%Daten für folgende Punktetabelle

\renewcommand{\aeins}{  3  }

\renewcommand{\azwei}{  3   }

\renewcommand{\adrei}{  6  }

\renewcommand{\avier}{  3  }

\renewcommand{\afuenf}{  3  }

\renewcommand{\asechs}{  5  }

\renewcommand{\asieben}{  13  }

\renewcommand{\aacht}{  4  }

\renewcommand{\aneun}{  3  }

\renewcommand{\azehn}{  7  }

\renewcommand{\aelf}{  6  }

\renewcommand{\azwoelf}{  3  }

\renewcommand{\adreizehn}{  6   }

\renewcommand{\avierzehn}{  65  }

\renewcommand{\afuenfzehn}{    }

\renewcommand{\asechzehn}{    }

\renewcommand{\asiebzehn}{    }

\renewcommand{\aachtzehn}{    }

\renewcommand{\aneunzehn}{    }

\renewcommand{\azwanzig}{    }

\renewcommand{\aeinundzwanzig}{    }

\renewcommand{\azweiundzwanzig}{    }

\renewcommand{\adreiundzwanzig}{    }

\renewcommand{\avierundzwanzig}{    }

\renewcommand{\afuenfundzwanzig}{    }

\renewcommand{\asechsundzwanzig}{    }

\punktetabelledreizehn

\klausurnote

\newpage

\setcounter{section}{0}

\inputaufgabegibtloesung
{3}
{

Definiere die folgenden
\zusatzklammer {kursiv gedruckten} {} {} Begriffe.
\aufzaehlungsechs{Ein
\stichwort {Hausdorff-Raum} {}
$X$.

}{Eine
\stichwort {differenzierbare Abbildung} {}
\maabbdisp {\varphi} { L } { M
} {}
zwischen zwei differenzierbaren Mannigfaltigkeiten
\mathkor {} {L} {und} {M} {.}

}{Das
\stichwort {Tangentialbündel} {}
zu einer differenzierbaren Mannigfaltigkeit $M$.

}{Die
\stichwort {zweite Fundamentalmatrix} {}
zu einer lokalen Parametrisierung
\maabbdisp {} {V} {U \subseteq Y
} {}
\zusatzklammer {mit

\mavergleichskette
{\vergleichskette
{ V
}
{ \subseteq }{ \R^{n-1}
}
{  }{
}
{    }{
}
{   }{
}
}
{}{}{}
offen} {} {}
einer differenzierbaren Hyperfläche

\mavergleichskette
{\vergleichskette
{ Y
}
{ \subseteq }{ \R^n
}
{  }{
}
{    }{
}
{   }{
}
}
{}{}{.}

}{Die \stichwort {äußere Ableitung} {} zu einer stetig differenzierbaren Differentialform
\mathl{\omega \in
{ \mathcal E }^{ k } (  M   )}{} auf einer differenzierbaren Mannigfaltigkeit $M$.

}{Eine \stichwort {der Überdeckung untergeordnete Partition der Eins} {,} wobei

\mavergleichskette
{\vergleichskette
{  X
}
{ = }{ \bigcup_{i \in I} W_i
}
{  }{
}
{    }{
}
{   }{
}
}
{}{}{}
eine
\definitionsverweis {offene Überdeckung}{}{}
eines
\definitionsverweis {topologischen Raumes}{}{}
$X$ ist.
}

}
{} {}

\inputaufgabegibtloesung
{3}
{

Formuliere die folgenden Sätze.
\aufzaehlungdrei{Der
\stichwort {Satz über die algebraische Eigenschaft der Weingartenabbildung} {.}}{Der Satz über Orientierbarkeit und positive Volumenform.}{Der
\stichwort {Brouwersche Fixpunktsatz} {.}}

}
{} {}

\inputaufgabegibtloesung
{6}
{

Bestimme die
\definitionsverweis {Krümmung}{}{}
in jedem Punkt des durch

\mavergleichskettedisp
{\vergleichskette
{ h(x,y)
}
{ =} { x^2+y^2
}
{ =} { 1
}
{ } {
}
{ } {
}
}
{}{}{}
implizit gegebenen Einheitskreises mit
[[Ebene differenzierbare Kurve/Faser/Krümmung/Fakt|Lemma 3.11 (Differentialgeometrie (Osnabrück 2023))]]
und dem Gradientenfeld zu $h$.

}
{} {}

\inputaufgabegibtloesung
{3}
{

Es sei $X$ ein
\definitionsverweis {Hausdorffraum}{}{.}
Zeige, dass die
\definitionsverweis {Diagonale}{}{}

\mavergleichskettedisp
{\vergleichskette
{ \triangle
}
{ =} { \left\{ (x,y) \in X \times X  \mid   x = y        \right\}
}
{ } {
}
{ } {
}
{ } {
}
}
{}{}{}
eine
\definitionsverweis {abgeschlossene Teilmenge}{}{}
im
\definitionsverweis {Produktraum}{}{}

\mathl{X \times X}{} ist.

}
{} {}

\inputaufgabegibtloesung
{3}
{

Zeige, dass die drei eindimensionalen Mannigfaltigkeiten
\mathlistdisp {S^1 = \left\{ (x,y) \in \R^2  \mid   \betrag { (x,y) } = 1         \right\}} {} {\R} {und} {\R \setminus \{0\}} {}
paarweise nicht homöomorph sind.

}
{} {}

\inputaufgabegibtloesung
{5}
{

Es sei

\mavergleichskette
{\vergleichskette
{ G
}
{ \subseteq }{ \R^m
}
{  }{
}
{    }{
}
{   }{
}
}
{}{}{}
\definitionsverweis {offen}{}{}
und sei

\mavergleichskette
{\vergleichskette
{ M
}
{ \subseteq }{ G
}
{  }{
}
{    }{
}
{   }{
}
}
{}{}{}
eine
$n$-\definitionsverweis {dimensionale}{}{}
\definitionsverweis {abgeschlossene Untermannigfaltigkeit}{}{,}
die
\definitionsverweis {orientiert}{}{}
und mit der induzierten riemannschen Struktur und der
\definitionsverweis {kanonischen Volumenform}{}{}
$\omega$ versehen sei. Es sei

\mavergleichskette
{\vergleichskette
{ W
}
{ \subseteq }{ \R^n
}
{  }{
}
{    }{
}
{   }{
}
}
{}{}{}
offen und es sei
\maabbdisp {\varphi} { W } { U
} {}
ein
\definitionsverweis {Diffeomorphismus}{}{}
mit der offenen Menge

\mavergleichskette
{\vergleichskette
{ U
}
{ \subseteq }{ M
}
{  }{
}
{    }{
}
{   }{
}
}
{}{}{.}
Zeige, dass auf $W$ die Formel

\mavergleichskettealignhandlinks
{\vergleichskettealignhandlinks
{\varphi^* \left(  \omega \vert_U  \right)
}
{ =} { \left( \det  \left( \left\langle  \begin{pmatrix}  { \frac{   \partial  \varphi_1   }{  \partial   x_i   } }  \\ \vdots \\  { \frac{   \partial  \varphi_m   }{  \partial   x_i   } }   \end{pmatrix}  ,  \begin{pmatrix}  { \frac{   \partial  \varphi_1   }{  \partial   x_j   } }  \\ \vdots \\  { \frac{   \partial  \varphi_m   }{  \partial   x_j   } }   \end{pmatrix}   \right\rangle \right)_{1 \leq i,j \leq n} \right)^{1/2}  dx_1 \wedge \ldots \wedge dx_n
}
{ =} { \left( \det  \left(  { \frac{   \partial  \varphi_1   }{  \partial   x_i   } } { \frac{   \partial  \varphi_1   }{  \partial   x_j   } }  + \cdots +  { \frac{   \partial  \varphi_m   }{  \partial   x_i   } } { \frac{   \partial  \varphi_m   }{  \partial   x_j   } }  \right)_{1 \leq i,j \leq n} \right)^{1/2}  dx_1 \wedge \ldots \wedge dx_n
}
{ } {
}
{ } {
}
}
{}
{}{}
gilt.

}
{} {}

\inputaufgabegibtloesung
{13 (2+3+2+2+2+2)}
{

Wir betrachten den

\mavergleichskette
{\vergleichskette
{ \R^6
}
{ = }{ \left(  \R^2  \right)^3
}
{  }{
}
{    }{
}
{   }{
}
}
{}{}{}
als Menge aller
\zusatzklammer {auch entarteter} {} {}
Dreiecke, indem wir ein Dreieck mit den
\zusatzklammer {geordneten} {} {}
Eckpunkten

\mavergleichskette
{\vergleichskette
{ P_1
}
{ = }{ \begin{pmatrix}  x_1  \\ y_1    \end{pmatrix}
}
{  }{
}
{    }{
}
{   }{
}
}
{}{}{,}

\mavergleichskette
{\vergleichskette
{ P_2
}
{ = }{ \begin{pmatrix}  x_2  \\ y_2    \end{pmatrix}
}
{  }{
}
{    }{
}
{   }{
}
}
{}{}{}
und

\mavergleichskette
{\vergleichskette
{ P_3
}
{ = }{ \begin{pmatrix}  x_3  \\ y_3    \end{pmatrix}
}
{  }{
}
{    }{
}
{   }{
}
}
{}{}{,}
mit dem Koordinatentupel

\mathdisp {(x_1,y_1,x_2,y_2,x_3,y_3)} {  }

identifizieren.
\aufzaehlungsechs{Zeige, dass die Menge der Dreiecke, bei denen zwei Eckpunkte zusammenfallen, eine abgeschlossene Teilmenge des $\R^6$ ist
\zusatzklammer {das Komplement davon ist somit eine offene Menge in $\R^6$, die wir $U$ nennen} {} {.}
}{Zeige, dass die Menge der Dreiecke, bei denen alle drei Eckpunkte auf einer Geraden liegen, eine abgeschlossene Teilmenge des $\R^6$ ist
\zusatzklammer {das Komplement davon, das aus allen nichtentarteten Dreiecken besteht, ist somit eine offene Menge in $\R^6$, die wir $V$ nennen} {} {.}
}{Erstelle eine Funktionsvorschrift, die die Abbildung
\maabbdisp {F} {\R^6} {\R
} {}
beschreibt, die einem Dreieck seinen Umfang zuordnet.
}{Zeige, dass die Funktion $F$ aus Teil (3) auf der Menge $U$ stetig differenzierbar ist.
}{Berechne die partielle Ableitung von $F$ nach $x_1$ auf $U$.
}{Zeige, dass die Menge der nichtentarteten Dreiecke mit Umfang $1$ eine
\definitionsverweis {abgeschlossene Untermannigfaltigkeit}{}{}
von $V$ bildet. Was ist die Dimension?
}

}
{} {}

\inputaufgabegibtloesung
{4}
{

Es seien
\mathkor {} {r} {und} {R} {}
positive reelle Zahlen mit

\mavergleichskette
{\vergleichskette
{ 0
}
{ < }{ r
}
{ < }{ R
}
{    }{
}
{   }{
}
}
{}{}{,}
wir betrachten den
\zusatzklammer {eingebetteten} {} {}
Torus

\mavergleichskettedisp
{\vergleichskette
{ T
}
{ =} { \left\{ (x,y,z) \in \R^3  \mid   \left(  \sqrt{x^2+y^2} -R  \right)^2 +z^2  = r^2         \right\}
}
{ \subseteq} { \R^3
}
{ } {
}
{ } {
}
}
{}{}{}
mit der
\definitionsverweis {induzierten riemannschen Struktur}{}{.}
Zeige, dass zu jedem

\mavergleichskette
{\vergleichskette
{ \alpha
}
{ \in }{ [0,2 \pi [
}
{  }{
}
{    }{
}
{   }{
}
}
{}{}{}
die Abbildung
\maabbeledisp {\Psi_\alpha} {\R^3} {\R^3
} { \left(  x   , \,   y  , \,  z                 \right) } { \left(  x \cos \alpha - y \sin  \alpha   , \,   x \sin \alpha + y \cos  \alpha  , \,  z                 \right)
} {,}
eine
\definitionsverweis {Isometrie}{}{}
von $T$ in sich induziert.

}
{} {}

\inputaufgabe
{3}
{

Beschreibe ein Modell für die hyperbolische Fläche.

}
{} {}

\inputaufgabegibtloesung
{7}
{

Beweise das Theorema egregium.

}
{} {}

\inputaufgabegibtloesung
{6 (1+2+2+1)}
{

Wir betrachten die
\definitionsverweis {Differentialform}{}{}

\mavergleichskettedisp
{\vergleichskette
{ \omega
}
{ =} { ydx+zdy +xdz
}
{ } {
}
{ } {
}
{ } {
}
}
{}{}{}
auf dem $\R^3$ und die Abbildung
\maabbeledisp {\varphi} {\R^2} {\R^3
} { \left(  u   , \,  v                   \right) } { \left( u^2  , \,  v^2  , \,  uv                 \right)
} {.}
\aufzaehlungvier{Berechne die äußere Ableitung von $\omega$.
}{Berechne den Rückzug $\varphi^* \omega$  von $\omega$ unter $\varphi$.
}{Berechne die äußere Ableitung von $\varphi^* \omega$ auf $\R^2$.
}{Berechne den Rückzug $\varphi^* d \omega$  von $d\omega$ unter $\varphi$ unabhängig von (3).
}

}
{} {}

\inputaufgabegibtloesung
{3}
{

Es sei $M$ eine
\definitionsverweis {differenzierbare Mannigfaltigkeit}{}{}
und $\omega$ eine
\definitionsverweis {exakte}{}{} \definitionsverweis {Differentialform}{}{}
auf $M$. Es sei $S$ eine
\definitionsverweis {kompakte}{}{}
\definitionsverweis {orientierte}{}{}
\definitionsverweis {differenzierbare Mannigfaltigkeit}{}{}
\zusatzklammer {ohne Rand} {} {}
mit
\definitionsverweis {abzählbarer Basis der Topologie}{}{}
und es sei
\maabbdisp {\varphi} { S } { M
} {}
eine
\definitionsverweis {stetig differenzierbare}{}{}
Abbildung. Zeige

\mavergleichskettedisp
{\vergleichskette
{
\int_{  S  }   \varphi^*\omega
}
{ =} { 0
}
{ } {
}
{ } {
}
{ } {
}
}
{}{}{.}

}
{} {}

\inputaufgabegibtloesung
{6}
{

Beweise den Satz über die Retraktionen auf den Rand auf Mannigfaltigkeiten.

}
{} {}

 [[Kategorie:Latexseite]]
